\documentclass[a4paper,12pt]{article}
\usepackage[utf8]{inputenc}
\usepackage[T1]{fontenc}
\usepackage[ngerman]{babel}
\usepackage{amsmath,amssymb}
\usepackage{booktabs}
\usepackage{enumitem}
\usepackage[margin=2.5cm]{geometry}

\begin{document}

\begin{center}
{\LARGE\bfseries \"Ubungsblatt 2 -- L\"osungen}\\[0.5em]
{\large VC Logik}
\end{center}

\bigskip

%% ============================================================
%% AUFGABE 1
%% ============================================================
\subsection*{Aufgabe 1 \textnormal{(2 Punkte)} -- Kaffeewelt}

Wissensbasis \"uber $V = \{G_{1,1},\, G_{1,2},\, K_{2,1},\, K_{2,2},\, K_{1,3}\}$:
\begin{align*}
\text{KB} &\equiv \neg G_{1,1} \wedge G_{1,2} \wedge {} \\
&\quad (K_{2,1} \to G_{1,1}) \wedge (K_{2,2} \to G_{1,2}) \wedge (K_{1,3} \to G_{1,2}) \wedge {} \\
&\quad (G_{1,1} \to K_{2,1}) \wedge (G_{1,2} \to K_{2,2} \vee K_{1,3})
\end{align*}

Die KB besteht aus 7~Teilen, verbunden mit $\wedge$ ("`und"').
Alle 7 m\"ussen gleichzeitig wahr sein:

\begin{enumerate}
\item $\neg G_{1,1}$: Kein Geruch in $(1{,}1)$ $\;\Rightarrow\;$ erzwingt $G_{1,1} = 0$
\item $G_{1,2}$: Geruch in $(1{,}2)$ $\;\Rightarrow\;$ erzwingt $G_{1,2} = 1$
\item $K_{2,1} \to G_{1,1}$: Tasse in $(2{,}1)$ $\Rightarrow$ Geruch in $(1{,}1)$.\\
      Da $G_{1,1}=0$: $K_{2,1} \to 0$ ist nur wahr, wenn $K_{2,1}=0$
      $\;\Rightarrow\;$ erzwingt $K_{2,1} = 0$
\item $K_{2,2} \to G_{1,2}$: Tasse in $(2{,}2)$ $\Rightarrow$ Geruch in $(1{,}2)$.\\
      Da $G_{1,2}=1$: $K_{2,2} \to 1 = 1$ f\"ur jedes $K_{2,2}$
      $\;\Rightarrow\;$ keine Einschr\"ankung
\item $K_{1,3} \to G_{1,2}$: Tasse in $(1{,}3)$ $\Rightarrow$ Geruch in $(1{,}2)$.\\
      Analog zu Teil~4: immer wahr $\;\Rightarrow\;$ keine Einschr\"ankung
\item $G_{1,1} \to K_{2,1}$: Geruch in $(1{,}1)$ $\Rightarrow$ Tasse in $(2{,}1)$.\\
      Da $G_{1,1}=0$: $0 \to K_{2,1} = 1$ (immer wahr)
      $\;\Rightarrow\;$ keine Einschr\"ankung
\item $G_{1,2} \to K_{2,2} \vee K_{1,3}$: Geruch in $(1{,}2)$ $\Rightarrow$
      mind.\ eine Tasse in $(2{,}2)$ oder $(1{,}3)$.\\
      Da $G_{1,2}=1$: $1 \to (K_{2,2} \vee K_{1,3})$ ist nur wahr,
      wenn $K_{2,2} \vee K_{1,3} = 1$
      $\;\Rightarrow\;$ erzwingt $K_{2,2} = 1$ oder $K_{1,3} = 1$
\end{enumerate}

\textbf{Zusammenfassung:} KB ist wahr gdw.\
$G_{1,1}=0$, $G_{1,2}=1$, $K_{2,1}=0$ und $K_{2,2} \vee K_{1,3} = 1$.

\newpage

\noindent\fbox{\parbox{\dimexpr\textwidth-2\fboxsep-2\fboxrule}{%
\small\textbf{Cheat-Sheet:}\;
$G_{1,1}\!=\!0$,\;
$G_{1,2}\!=\!1$,\;
$K_{2,1}\!=\!0$,\;
$K_{2,2}$\,/\,$K_{1,3}$ frei aber $K_{2,2} \vee K_{1,3}\!=\!1$
\;$\Rightarrow$\;
$(K_{2,2}, K_{1,3}) \in \{(0,1),\,(1,0),\,(1,1)\}$
}}

\bigskip
Es gibt $2^5 = 32$ Interpretationen.
Alle 32 mit Auswertung von~KB
(Modelle \textbf{fett}, Welten aus der Vorlesung in der letzten Spalte):

\begin{center}
\small
\renewcommand{\arraystretch}{1.05}
\begin{tabular}{r|ccccc|c|c}
\toprule
\# & $G_{1,1}$ & $G_{1,2}$ & $K_{2,1}$ & $K_{2,2}$ & $K_{1,3}$
   & KB & Welt \\
\midrule
1  & 0 & 0 & 0 & 0 & 0 & 0 & $W_1$ \\
2  & 0 & 0 & 0 & 0 & 1 & 0 & \\
3  & 0 & 0 & 0 & 1 & 0 & 0 & \\
4  & 0 & 0 & 0 & 1 & 1 & 0 & \\
5  & 0 & 0 & 1 & 0 & 0 & 0 & \\
6  & 0 & 0 & 1 & 0 & 1 & 0 & \\
7  & 0 & 0 & 1 & 1 & 0 & 0 & \\
8  & 0 & 0 & 1 & 1 & 1 & 0 & \\
\midrule
9  & 0 & 1 & 0 & 0 & 0 & 0 & \\
10 & 0 & 1 & 0 & 0 & 1 & \textbf{1} & $W_2$ \\
11 & 0 & 1 & 0 & 1 & 0 & \textbf{1} & $W_3$ \\
12 & 0 & 1 & 0 & 1 & 1 & \textbf{1} & $W_4$ \\
13 & 0 & 1 & 1 & 0 & 0 & 0 & \\
14 & 0 & 1 & 1 & 0 & 1 & 0 & \\
15 & 0 & 1 & 1 & 1 & 0 & 0 & \\
16 & 0 & 1 & 1 & 1 & 1 & 0 & \\
\midrule
17 & 1 & 0 & 0 & 0 & 0 & 0 & \\
18 & 1 & 0 & 0 & 0 & 1 & 0 & \\
19 & 1 & 0 & 0 & 1 & 0 & 0 & \\
20 & 1 & 0 & 0 & 1 & 1 & 0 & \\
21 & 1 & 0 & 1 & 0 & 0 & 0 & $W_5$ \\
22 & 1 & 0 & 1 & 0 & 1 & 0 & \\
23 & 1 & 0 & 1 & 1 & 0 & 0 & \\
24 & 1 & 0 & 1 & 1 & 1 & 0 & \\
\midrule
25 & 1 & 1 & 0 & 0 & 0 & 0 & \\
26 & 1 & 1 & 0 & 0 & 1 & 0 & \\
27 & 1 & 1 & 0 & 1 & 0 & 0 & \\
28 & 1 & 1 & 0 & 1 & 1 & 0 & \\
29 & 1 & 1 & 1 & 0 & 0 & 0 & \\
30 & 1 & 1 & 1 & 0 & 1 & 0 & $W_6$ \\
31 & 1 & 1 & 1 & 1 & 0 & 0 & $W_7$ \\
32 & 1 & 1 & 1 & 1 & 1 & 0 & $W_8$ \\
\bottomrule
\end{tabular}
\end{center}

\medskip
\textbf{Begr\"undung der Gruppen:}
\begin{itemize}[nosep]
\item Nr.\,1--8 ($G_{1,2}=0$): KB verlangt $G_{1,2}=1$ (Teil~2)
      $\Rightarrow$ sofort falsch.
\item Nr.\,9--16 ($G_{1,1}=0$, $G_{1,2}=1$): Beobachtungen erf\"ullt.
      \begin{itemize}[nosep]
      \item Nr.\,9: $K_{2,2}=K_{1,3}=0$ $\Rightarrow$ Teil~7 verletzt (kein Geruchsursprung).
      \item Nr.\,10--12: $K_{2,1}=0$ und $K_{2,2} \vee K_{1,3} = 1$
            $\Rightarrow$ alle 7~Teile wahr $\Rightarrow$ \textbf{Modelle}.
      \item Nr.\,13--16: $K_{2,1}=1$ $\Rightarrow$ Teil~3 verletzt ($1 \to 0 = 0$).
      \end{itemize}
\item Nr.\,17--32 ($G_{1,1}=1$): KB verlangt $\neg G_{1,1}$ (Teil~1)
      $\Rightarrow$ sofort falsch.
\end{itemize}

\bigskip
\textbf{Verh\"altnis zu den 8~Welten:}

Die 8~Welten $W_1$--$W_8$ aus der Vorlesung entsprechen den $2^3 = 8$
Kombinationen der Tassenvariablen ($K_{2,1}$, $K_{2,2}$, $K_{1,3}$),
wobei die Geruchswerte durch die Regeln eindeutig bestimmt sind:
$G_{1,1} = K_{2,1}$ und $G_{1,2} = K_{2,2} \vee K_{1,3}$.
Die restlichen $32-8 = 24$ Interpretationen verletzen diese
Konsistenz und machen~KB automatisch falsch.

\begin{center}
\renewcommand{\arraystretch}{1.2}
\begin{tabular}{c|ccc|cc|c|l}
\toprule
Welt & $K_{2,1}$ & $K_{2,2}$ & $K_{1,3}$
     & $G_{1,1}$ & $G_{1,2}$ & Modell? & Grund \\
\midrule
$W_1$ & 0 & 0 & 0 & 0 & 0 & Nein & $G_{1,2} \neq 1$ (kein Geruchsursprung) \\
$W_2$ & 0 & 0 & 1 & 0 & 1 & \textbf{Ja} & Tasse nur in $(1{,}3)$ \\
$W_3$ & 0 & 1 & 0 & 0 & 1 & \textbf{Ja} & Tasse nur in $(2{,}2)$ \\
$W_4$ & 0 & 1 & 1 & 0 & 1 & \textbf{Ja} & Tassen in $(2{,}2)$ und $(1{,}3)$ \\
$W_5$ & 1 & 0 & 0 & 1 & 0 & Nein & $G_{1,1} \neq 0$ (Tasse in $(2{,}1)$) \\
$W_6$ & 1 & 0 & 1 & 1 & 1 & Nein & $G_{1,1} \neq 0$ \\
$W_7$ & 1 & 1 & 0 & 1 & 1 & Nein & $G_{1,1} \neq 0$ \\
$W_8$ & 1 & 1 & 1 & 1 & 1 & Nein & $G_{1,1} \neq 0$ \\
\bottomrule
\end{tabular}
\end{center}

\textbf{Ergebnis:} 3~Modelle ($W_2$, $W_3$, $W_4$).
In $(2{,}1)$ steht sicher keine Tasse; in $(2{,}2)$ oder $(1{,}3)$
(oder beiden) muss mindestens eine stehen.

\newpage

%% ============================================================
%% AUFGABE 2
%% ============================================================
\subsection*{Aufgabe 2 \textnormal{(3 Punkte)} -- G\"ultigkeit,
Erf\"ullbarkeit, \"Aquivalenz, log.\ Konsequenz}

Formeln \"uber $V = \{P, Q\}$:
\;(a)\,$P$,
\;(b)\,$Q$,
\;(c)\,$P \wedge Q$,
\;(d)\,$\neg P \vee \neg Q$,
\;(e)\,$\neg(\neg P \vee \neg Q)$.

\medskip
Wahrheitstabelle ($2^2 = 4$ Zeilen, weil 2~Variablen):

\begin{center}
\renewcommand{\arraystretch}{1.25}
\begin{tabular}{cc|ccccc}
\toprule
$P$ & $Q$
  & (a)\;$P$
  & (b)\;$Q$
  & (c)\;$P \wedge Q$
  & (d)\;$\neg P \vee \neg Q$
  & (e)\;$\neg(\neg P \vee \neg Q)$ \\
\midrule
0 & 0 & 0 & 0 & 0 & 1 & 0 \\
0 & 1 & 0 & 1 & 0 & 1 & 0 \\
1 & 0 & 1 & 0 & 0 & 1 & 0 \\
1 & 1 & 1 & 1 & 1 & 0 & 1 \\
\bottomrule
\end{tabular}
\end{center}

\bigskip
\noindent\fbox{\parbox{\dimexpr\textwidth-2\fboxsep-2\fboxrule}{%
\renewcommand{\arraystretch}{1.2}
\begin{tabular}{@{}l@{\;\;}l@{}}
\textbf{G\"ultigkeit:} & Keine (keine Spalte hat nur Einsen). \\
\textbf{Erf\"ullbarkeit:} & Alle f\"unf (jede Spalte hat mind.\ eine Eins). \\
\textbf{\"Aquivalenz:} & $(c) \equiv (e)$
  \;(De~Morgan: $P \wedge Q \equiv \neg(\neg P \vee \neg Q)$). \\
\end{tabular}
}}

\bigskip
\textbf{Logische Konsequenz} ($F_1 \models F_2$: \"uberall wo $F_1 = 1$,
muss auch $F_2 = 1$ sein):

\begin{center}
\renewcommand{\arraystretch}{1.15}
\begin{tabular}{l@{\quad}l}
\toprule
\textbf{Gilt} & \textbf{Begr\"undung} \\
\midrule
$(c) \models (a)$ & (c) nur in Z.\,4 wahr; dort (a)$=1$ \checkmark \\
$(c) \models (b)$ & (c) nur in Z.\,4 wahr; dort (b)$=1$ \checkmark \\
$(e) \models (a)$, $(e) \models (b)$ & (e) hat gleiche Spalte wie (c) \\
\midrule
\textbf{Gilt nicht} & \textbf{Gegenbeispiel} \\
\midrule
$(a) \not\models (c)$ & Z.\,3: (a)$=1$, (c)$=0$ \\
$(a) \not\models (b)$ & Z.\,3: (a)$=1$, (b)$=0$ \\
$(b) \not\models (a)$ & Z.\,2: (b)$=1$, (a)$=0$ \\
$(d) \not\models (a)$ & Z.\,1: (d)$=1$, (a)$=0$ \\
\bottomrule
\end{tabular}
\end{center}
Keine weiteren Konsequenzen.

\newpage

%% ============================================================
%% AUFGABE 3
%% ============================================================
\subsection*{Aufgabe 3 \textnormal{(2 Punkte)} -- Erf\"ullbarkeitstester}

Ein Erf\"ullbarkeitstester ist ein Automat, der eine Formel bekommt
und nur mit "`erf\"ullbar"' oder "`unerf\"ullbar"' antwortet.
Um andere Fragen damit zu beantworten, formuliere ich sie so um,
dass ich das Ergebnis aus seiner Antwort ablesen kann.

\medskip
\noindent\fbox{\parbox{\dimexpr\textwidth-2\fboxsep-2\fboxrule}{%
\renewcommand{\arraystretch}{1.3}
\begin{tabular}{@{}l@{\;\;}l@{\;\;}l@{}}
\textbf{Frage} & \textbf{Dem Tester geben}
  & \textbf{Wenn unerf\"ullbar} \\
\hline
Ist $F$ g\"ultig / Tautologie? & $\neg F$
  & $\Rightarrow$ ja, g\"ultig \\
Ist $F$ eine Kontradiktion? & $F$ direkt
  & $\Rightarrow$ ja, Kontradiktion \\
Folgt $F_2$ aus $F_1$? & $F_1 \wedge \neg F_2$
  & $\Rightarrow$ ja, es folgt \\
\end{tabular}
}}

\bigskip
Damit ergeben sich die Umformulierungen und Ergebnisse:

\begin{center}
\renewcommand{\arraystretch}{1.4}
\begin{tabular}{c|l|l|l}
\toprule
 & \textbf{Frage} & \textbf{Dem Tester geben}
 & \textbf{Ergebnis} \\
\midrule
(a) & G\"ultig?
  & $\neg\bigl((P \vee (\neg P \to Q)) \leftrightarrow (P \vee Q)\bigr)$
  & unerf. $\Rightarrow$ \textbf{g\"ultig} \\
(b) & Folgt $P \!\to\! Q$ aus $\neg Q \!\to\! \neg P$?
  & $(\neg Q \to \neg P) \wedge \neg(P \to Q)$
  & unerf. $\Rightarrow$ \textbf{ja} \\
(c) & Kontradiktion?
  & $P \leftrightarrow (Q \wedge P)$ direkt
  & erf\"ullbar $\Rightarrow$ \textbf{nein} \\
(d) & Tautologie?
  & $\neg\bigl(P \leftrightarrow (P \vee \bot)\bigr)$
  & unerf. $\Rightarrow$ \textbf{ja} \\
\bottomrule
\end{tabular}
\end{center}

\medskip
\textbf{Anmerkungen:}
\begin{itemize}[nosep]
\item (c): $\wedge$ bindet st\"arker als $\leftrightarrow$,
      daher $P \leftrightarrow (Q \wedge P)$.
\item (d): $\bot = 0$ (immer falsch), also
      $P \vee \bot = P$ und die Formel wird zu $P \leftrightarrow P$.
\end{itemize}

\medskip
\textit{Flei\ss aufgaben (Schritt f\"ur Schritt vereinfachen):}

\medskip
\textbf{(a)} Die Formel ist
$(P \vee (\neg P \to Q)) \leftrightarrow (P \vee Q)$.

Den mittleren Teil vereinfachen mit der Regel
$A \to B \equiv \neg A \vee B$:
\[
\neg P \to Q \;=\; \neg(\neg P) \vee Q \;=\; P \vee Q
\]
Einsetzen:
$(P \vee \underbrace{(\neg P \to Q)}_{= P \vee Q})
\leftrightarrow (P \vee Q)
\;=\; (P \vee P \vee Q) \leftrightarrow (P \vee Q)$

$P \vee P \vee Q$ ist einfach $P \vee Q$, also steht da:
$(P \vee Q) \leftrightarrow (P \vee Q)$

Das hat die Form $X \leftrightarrow X$ -- immer wahr
$\Rightarrow$ \textbf{g\"ultig}.

\medskip
\textbf{(b)} Die Frage ist: Folgt $P \to Q$ aus $\neg Q \to \neg P$?

$\neg Q \to \neg P$ umschreiben mit $A \to B \equiv \neg A \vee B$:
\[
\neg Q \to \neg P \;=\; \neg(\neg Q) \vee \neg P \;=\; Q \vee \neg P
  \;=\; \neg P \vee Q
\]
Und $P \to Q$ genauso:
\[
P \to Q \;=\; \neg P \vee Q
\]
Beide Formeln sind identisch ($\neg P \vee Q$), also \"aquivalent.
Wenn die Pr\"amisse wahr ist, ist die Folgerung automatisch auch wahr
$\Rightarrow$ \textbf{ja, es folgt}.

\medskip
\textbf{(c)} Die Formel ist $P \leftrightarrow (Q \wedge P)$.

Ich suche eine Belegung, die die Formel wahr macht:

$P = 0$: \;$0 \leftrightarrow (Q \wedge 0)
= 0 \leftrightarrow 0 = 1$ \checkmark

Schon mit $P=0$ (egal welches $Q$) ist die Formel wahr.
Also erf\"ullbar $\Rightarrow$ \textbf{keine Kontradiktion}.

\medskip
\textbf{(d)} Die Formel ist $P \leftrightarrow (P \vee \bot)$.

$\bot$ ist immer falsch (= 0). Vereinfachen:
\[
P \vee \bot \;=\; P \vee 0 \;=\; P
\]
Also wird die Formel zu $P \leftrightarrow P$.

Das hat wieder die Form $X \leftrightarrow X$ -- immer wahr
$\Rightarrow$ \textbf{Tautologie}.

\end{document}
